\section{Gradient flow and small flow-time expansion}
\label{sec:introGF}
\subsection{The gradient-flow formalism}
In the following, we work in $d$-dimensional Euclidean spacetime with $d=4-2\epsilon$.
The gradient flow extends the $d$-dimensional QCD to a $(d+1)$-dimensional field theory by introducing an additional coordinate $t$ (not to be confused with the time coordinate in Minkowski spacetime) of mass dimension $-2$,  which is called the flow time \cite{Luscher:2010iy, Luscher:2011bx}. In the gradient-flow formalism, the flow time dependent gluon field $B_{\mu}^a(x, t)$ and quark field $\chi^i_{\alpha}(x, t)$ 
are determined by the flow equations \cite{Luscher:2010iy, Luscher:2011bx,Luscher:2013cpa},
\begin{eqnarray}
\label{eq:flowequations}
\partial_t B_{\mu} &=& \mathcal{D}_{\nu}G_{\nu\mu} +\kappa \mathcal{D}_{\mu}\partial_{\nu}B_{\nu},\\
\partial\chi &= & \mathcal{D}_{\mu}\mathcal{D}_{\mu}\chi -\kappa (\partial_{\mu}B_{\mu})\chi,\\
\partial\bar{\chi} &=& \bar{\chi}\overleftarrow{\mathcal{D}}_{\mu}\overleftarrow{\mathcal{D}}_{\mu} +\kappa\bar{\chi}\partial_{\mu}B_{\mu},
\end{eqnarray}
where  $\kappa$ is a gauge parameter that drops out in physical observables, and
\begin{eqnarray}
G_{\mu\nu} &= &\partial_{\mu}B_{\nu} -\partial_{\nu}B_{\mu} + \left[B_{\mu}, B_{\nu}\right],\\
\mathcal{D}_{\mu} &=&\partial_{\mu} + \left[B_{\mu}, .\right],
\end{eqnarray}
with the boundary conditions
\begin{eqnarray}
\label{eq:boundary}
B_{\mu}(x; t=0) &=& gA_{\mu}(x),\\
\chi(x; t=0) &=& \psi(x),\\
\bar{\chi}(x; t=0) &=& \bar{\psi}(x).
\end{eqnarray}


In perturbative gradient flow, one extends the regular QCD Feynman rules by adding flow time dependent exponentials to the propagators, introducing  extra ``flow lines” and ``flow vertices” \cite{Luscher:2011bx}; see ref.~\cite{Artz:2019bpr} for more details about the gradient flow Feynman rules and higher order calculation techniques. In the following perturbative calculation, we  adopt the  convenient choice $\kappa=1$.

With renormalized physical parameters, the flowed field $B_{\mu}$ does not need further renormalization \cite{Luscher:2010iy,Luscher:2011bx}, but the flowed quark fields $\chi$ need to be  renormalized \cite{Luscher:2013cpa},
\begin{eqnarray}
\chi_{0}(x; t) = Z_{\chi}^{\frac{1}{2}} \chi(x; t), \, \, \bar{\chi}_0(x; t) = Z_{\chi}^{\frac{1}{2}} \bar{\chi}(x; t), \, \,  Z_{\chi} = 1- \frac{\alpha_s}{4\pi}C_F \frac{3}{\epsilon_{\rm UV}} + O(\alpha_s^2),
\end{eqnarray}
where $C_F = (N_c^2-1)/(2N_c)$ with $N_c=3$ the number of colors.
In order to avoid the complication due to the matching between $ Z_{\chi}$ in dimensional regularization and that in the lattice regularization, the ringed fermion fields $\mathring{\chi}(x,t)$, $\mathring{\bar{\chi}}(x,t)$ were introduced in ref.~\cite{Makino:2014taa},
\begin{eqnarray}
\mathring{\chi}(x,t) = \mathring{Z}_{\chi}^{\frac{1}{2}}\chi_{0}(x,t),\, \,  \, \mathring{\bar{\chi}}(x,t) = \mathring{Z}_{\chi}^{\frac{1}{2}}\bar{\chi}_{0}(x,t),
\end{eqnarray}
with
\begin{eqnarray}
 \mathring{Z}_{\chi}=\frac{-2N_c}{(4\pi)^2t^2
\langle \bar{\chi}_{0}(x,t)\overleftrightarrow{\slashed{\mathcal{D}}}\chi_{0}(x,t)\rangle|_{m=0}},
\end{eqnarray}
for one quark flavor, where $\overleftrightarrow{\mathcal{D}}^{\mu} = {\mathcal{D}}^{\mu}  - \overleftarrow{\mathcal{D}}^{\mu}$. The next-next-to-leading order (nnLO) result of $\mathring{Z}_{\chi}$ can be found in  ref.~\cite{Artz:2019bpr}. At next-to-leading order (NLO), the renormalization factor $ \mathring{Z}_{\chi}$ is given by \cite{Makino:2014taa}
\begin{eqnarray}
\label{eq:ringedfermionren}
\mathring{Z}_{\chi}(t, \mu) = (8\pi t)^{-\epsilon}\left\{1+  \frac{\alpha_s}{4\pi}C_F\left[\frac{3}{\epsilon_{\rm UV}} + 3\log\left(2\mu^2 te^{\gamma_E}\right)-\log(432)\right] \right\},
\end{eqnarray}
where  the factor $(8\pi t)^{-\epsilon}$ is introduced to balance the dimension of the ringed fermion field, which has dimension $3/2$, while the flowed fermion field has dimension $(d-1)/2$.




%%%%%%%%%%%%%%%%%%%%%%%%%%%%%%%%%%%%%%%%%%
\subsection{The small flow-time expansion}
\label{subsec:smallflowexp}
In this work, we consider the local current operators defined in eqs.~(\ref{eq:currentoperator}) and eq.~(\ref{eq:currentoperatorspin}), which are flowed at the flow time $t$. We use the same symbols for both the flowed and the un-flowed current operators as long as it does not lead to confusions. In the context of Wilson-line operators, we flow the external fields and the Wilson lines at the flow time $t$. Specifically, this means that all the gluon fields directly connecting to the ``heavy” auxiliary field lines are flowed at the flow time $t$, while the ``heavy” auxiliary fields themselves remain un-flowed. In the regime of small flow time, we employ a ``short-distance”  OPE  for the flowed current operators. Our objective is to derive matching coefficients between renormalized flowed current operators and the  $\overline{\rm MS}$ renormalized un-flowed current operators,
\begin{eqnarray}
\label{eq:smalltexpansion}
\mathcal{O}^{\rm R} (t)= c_{\mathcal{O}}(t, \mu)\mathcal{O}^{\overline{\rm MS}} (\mu) + O(t),
\end{eqnarray}
where we have neglected operator mixing. 


We work in the massless limit for quark fields and set all the external scales to be zero whenever feasible. 
This strategic choice allows us to leverage established techniques from HQET for matching calculations.  Specifically, we do not need to calculate the un-flowed operators because they yield scaleless integrals to all orders in perturbative calculations in dimensional regularization so that we can obtain the matching coefficients $c_{\mathcal{O}}(t, \mu)$ from the renormalized flowed operators simply by subtracting the IR poles.

Similarly to the renormalization of current operators, we can break down the matching coefficients $c_{\mathcal{O}}(t, \mu)$ for the local current operators into products of field matching coefficients and corresponding vertex matching coefficients. For the four local current operators $\bar{\psi} h_v$, $gF^{\mu\nu}_{\|\perp}h_v$, $gF^{\mu\nu}_{\perp\perp}h_v$, and $g\bar{h}_vF^{\mu\nu}_{\perp\perp}h_v$, the relationships between matching coefficients for the local current operators and those for the corresponding fields and vertices are given by
\begin{subequations}
\label{eq:operatorrelations}
\begin{eqnarray}
c_{\psi h_v} (t,\mu) &=& \mathring{\zeta}_{\psi}(t,\mu)\zeta_{h_v}^F(t,\mu)\zeta_{\psi h_v}(t,\mu),\\
c_{\|\perp} (t,\mu)&=& \zeta_{A}(t,\mu)\zeta_{h_v}^A(t,\mu)\zeta_{\|\perp}(t,\mu),\\
c_{\perp\perp}(t,\mu) &=& \zeta_{A}(t,\mu)\zeta_{h_v}^A(t,\mu)\zeta_{\perp\perp}(t,\mu),\\
c_F(t,\mu) &=& \zeta_{A}(t,\mu)(\zeta^F_{h_v}(t,\mu))^2\zeta_F(t,\mu).
\end{eqnarray}
\end{subequations}
Here $ \mathring{\zeta}_{\psi}(t,\mu)$, $\zeta_{h_v}^F (\zeta_{h_v}^A)$ and $\zeta_{A}(t,\mu)$ are the matching coefficients of the ringed fermion field $\mathring{\chi}$,  the ``heavy” auxiliary field $h_v$ in fundamental (adjoint) representation and the flowed gluon field $B$, respectively. Additionally, $\zeta_{\psi h_v}(t,\mu)$, $\zeta_{\|\perp}(t,\mu)$, $\zeta_{\perp\perp}(t,\mu)$ and $\zeta_F(t,\mu)$ denote matching coefficients stemming from vertex corrections for the corresponding flowed local current operators. We calculate these matching coefficients to one-loop order in section \ref{sec:oneloop}.

A key insight arising from our study in section \ref{sec:EFT} on the multiplicative renormalization of Wilson-line operators is that matching coefficients for Wilson-line operators in the transition from the gradient-flow scheme to the $\overline{\rm MS}$ scheme can be directly derived from matching coefficients for the corresponding local current operators. This occurs because in the one-dimensional auxiliary-field formalism, combinations of the renormalized local current operators at different spacetime points do not introduce additional UV divergences, apart from the subtracted linear divergence, and because of eq.~(\ref{eq:smalltexpansion}) that relates local operators in the small flow-time limit. Therefore, additional matching for these combined local current operators in different spacetime points is not needed in the small flow-time limit.




%%%%%%%%%%%%%%%%%%%%%%%%%%%%%%%%%%%%%%%%%%%%%%%%%%%%%%%%%
